% ======================================================================
% 语法着色测试文件 — 覆盖所有 TeX 内容类型
% 颜色注释格式:  ← 应显示为 XXX 色
% ======================================================================

% --- 0. 注释 (syntax_comment: 深色 #6A9955 / 浅色 #008000) -------------
% 本行及上面所有 % 开头的行都应显示为绿色
% 注释中的 \command 不应变黄，$E=mc^2$ 也不应变红
% 注释中的 \begin \end \usepackage 也不应变粉

% --- 1. 关键字命令 (tex_keyword: 深色 #C586C0 / 浅色 #AF00DB) ---------
% 以下命令应显示为品红/粉色：
\documentclass{article}     % ← \documentclass 是关键字命令 → tex_keyword 粉色
\usepackage{amsmath}        % ← \usepackage 是关键字命令 → tex_keyword 粉色
\usepackage{graphicx}       % ← \usepackage 是关键字命令 → tex_keyword 粉色

% --- 2. 函数命令 (tex_command: 深色 #DCDCAA / 浅色 #795E26) -----------
% 以下命令应显示为金色/棕色：
\section{函数命令测试}      % ← \section 是函数命令 → tex_command 金色
\subsection{小节标题}       % ← \subsection 是函数命令 → tex_command 金色
\textbf{粗体文字}           % ← \textbf 是函数命令 → tex_command 金色
\textit{斜体文字}           % ← \textit 是函数命令 → tex_command 金色
\ref{tab:data}              % ← \ref 是函数命令 → tex_command 金色
\cite{ref1}                 % ← \cite 是函数命令 → tex_command 金色
\label{sec:intro}           % ← \label 是函数命令 → tex_command 金色
\caption{表格标题}          % ← \caption 是函数命令 → tex_command 金色
\centering                  % ← \centering 是函数命令 → tex_command 金色

% --- 3. 带星号(*)的命令变体 --------------------------------------------
% \section* 整体(含*) → tex_command 金色
% \newcommand* 整体(含*) → tex_keyword 粉色
\section*{不带编号的节}     % ← \section* 含 * → tex_command 金色
\newcommand*{\mycmd}{foo}   % ← \newcommand* 含 * → tex_keyword 粉色（因为 \newcommand 是关键字）

% --- 4. 关键字命令 — 环境与定义 ----------------------------------------
% 以下都是 tex_keyword 粉色：
\begin{document}            % ← \begin 是关键字命令 → tex_keyword 粉色
\end{document}              % ← \end 是关键字命令 → tex_keyword 粉色

\begin{equation}            % ← \begin 是关键字命令 → tex_keyword 粉色
E = mc^{2}
\label{eq:einstein}         % ← \label 是函数命令 → tex_command 金色
\end{equation}              % ← \end 是关键字命令 → tex_keyword 粉色

\begin{table}[H]             % ← \begin 是关键字命令 → tex_keyword 粉色
\centering
\caption{测试数据}
\label{tab:data}
\end{table}                 % ← \end 是关键字命令 → tex_keyword 粉色

% --- 5. 文件包含类的关键字命令 -----------------------------------------
% 以下应显示为 tex_keyword 粉色：
\section{引言}

半导体器件设计与制造的核心问题之一是如何精确控制半导体内的杂质分布，以满足不同器件对电学特性的要求。杂质浓度及其分布直接影响pn结的击穿电压、势垒电容、载流子迁移率等关键参数，因此准确测量杂质分布是半导体材料与器件研究中的基本任务之一。

传统的杂质分布测量方法包括四探针法和霍耳效应逐次去层测量技术。这些方法通过逐层剥离样品并测量薄层电导或薄层霍耳电压，可以获得杂质浓度及其分布。然而，逐次去层测量操作繁琐、耗时长，且对样品具有破坏性。相比之下，电容-电压法（C--V法）利用pn结势垒电容随反向偏压变化的特性，通过测量不同直流偏压下的势垒电容，可以在不破坏器件的情况下快速获得单边突变pn结轻掺杂一侧的杂质浓度及其分布。该方法尤其适用于检测扩散工艺或离子注入工艺所引入的杂质分布。

本实验基于C--V法测量$\mathrm{p}^+\mathrm{n}$单边突变结轻掺杂一侧的杂质浓度分布。实验采用锁定放大器作为微弱信号检测的核心仪器，将pn结势垒电容转换为交流电压信号进行精确测量。通过对多个二极管样品的C--V特性进行系统测量，结合理论模型计算得到各二极管的杂质浓度分布$N(w)$，并通过$1/C^{2}$--$V_{\mathrm{R}}$曲线的直线段外推获得接触电势差$V_{\mathrm{D}}$。此外，实验还考察了锁定放大器在不同时间常量下抑制噪声的能力，以及二极管相位特性与反向偏压的关系。
      % ← \input 是关键字命令 → tex_keyword 粉色
\include{appendix}          % ← \include 是关键字命令 → tex_keyword 粉色
\bibliography{refs}         % ← \bibliography 是关键字命令 → tex_keyword 粉色
\bibliographystyle{plain}   % ← \bibliographystyle 是关键字命令 → tex_keyword 粉色

% --- 6. 定义类的关键字命令 ---------------------------------------------
% 以下应显示为 tex_keyword 粉色：
\newcommand{\foo}{bar}      % ← \newcommand 是关键字命令 → tex_keyword 粉色
\renewcommand{\foo}{baz}    % ← \renewcommand 是关键字命令 → tex_keyword 粉色
\newenvironment{myenv}{}{}  % ← \newenvironment 是关键字命令 → tex_keyword 粉色
\def\mydef{value}           % ← \def 是关键字命令 → tex_keyword 粉色
\let\oldcmd\newcmd          % ← \let 是关键字命令 → tex_keyword 粉色
\newif\ifmycond             % ← \newif 是关键字命令 → tex_keyword 粉色
\providecommand{\foo}{bar}  % ← \providecommand 是关键字命令 → tex_keyword 粉色

% --- 7. 数学模式 ($...$ / $$...$$: tex_math 深色 #CE9178 / 浅色 #A31515)
% 整个 $...$ 和 $$...$$ 范围内的内容应显示为橙色/红色
$V_{\rm out} = k C_{\rm nominal} + b$       % ← 整段 → tex_math 橙色
$R^2 = 0.9999$                                % ← 整段 → tex_math 橙色
$\phi = -171.93^\circ$                        % ← 整段 → tex_math 橙色
$\sigma \propto 1 / \sqrt{T_c}$               % ← 整段 → tex_math 橙色

% 行间公式（display math）—— 同样 tex_math 橙色：
$$N(w) = \frac{2}{q\epsilon\epsilon_0 A^2 \, d(1/C^2)/dV}$$

% --- 8. 分组括号 (tex_group: 深色 #D4D4D4 / 浅色 #000000) --------------
% { 和 } 应显示为灰色/黑色（标点色）
% 以下每一对 {} 中的括号：
%   {article}   ← 括号 → tex_group 灰色  内容 → editor_fg 正文色
%   {amsmath}   ← 括号 → tex_group 灰色
%   [H]         ← 括号 → tex_group 灰色（方括号也是 group）
% 下面这个命令被 { } 包围的参数括号都会用到 tex_group：
\section{分组括号测试：\textbf{内层粗体}}  % ← 外层{}和内层{} → tex_group 灰色

% --- 9. 特殊字符 (tex_special: 深色 #D4D4D4 / 浅色 #000000) ------------
% 以下特殊字符应显示为灰色/黑色（与分组括号同色）：
% 注意：这些字符必须出现在正文中，不能写在注释里！

\section{特殊字符正文测试}

% 在正文中使用特殊字符（通过 \texttt 或直接使用）：
% $  #  %  &  ~  _  ^  \  {  }  |
% 下面用 \verb 展示或实际嵌入：

美元符号可以直接使用: \$        % ← \$ 中的 $ → tex_special 灰色
井号: \#                         % ← \# 中的 # → tex_special 灰色
百分号: \%                       % ← \% 中的 % → tex_special 灰色
and符号: \&                      % ← \& 中的 & → tex_special 灰色
波浪线: \~{}                     % ← \~{} 中的 ~ → tex_special 灰色
下划线: \_                       % ← \_ 中的 _ → tex_special 灰色
帽子: \^{}                       % ← \^{} 中的 ^ → tex_special 灰色
反斜杠: \textbackslash           % ← 反斜杠 → tex_special 灰色
竖线在数学模式: $|x|$            % ← | 在数学模式中 → tex_special 灰色

% 直接出现在正文中的特殊字符（通过转义）：
% 实际上在 LaTeX 中，$ # % & ~ _ ^ \ { } 这些字符大多需要通过命令或
% 特定上下文才能出现。关键是它们在 QScintilla TeX lexer 中会被标记
% 为 "Special" 样式 → tex_special 色。

% --- 10. 混合场景测试 --------------------------------------------------
\section{混合场景}

根据公式 $V_{\rm out} \propto C_{\rm x}$，
% ↑ $...$ → tex_math 橙色，\rm → tex_command 金色
当 $C_0 \gg C_{\rm x}$ 时，锁定放大器输出与被测电容近似成正比。
其中 \textbf{灵敏度} 由 \textit{最小二乘法} 确定（见表~\ref{tab:calibration}）。
% ↑ \textbf → tex_command 金色，{} → tex_group 灰色

% 嵌套命令中的颜色优先级：
% \section{标题中的 \textbf{粗体} 和 $E=mc^2$}
%   ↑ section → tex_command    ↑ textbf → tex_command   ↑ 数学 → tex_math
%   ↑ { → tex_group            ↑ { → tex_group          ↑ { → tex_group

% --- 11. 引用与文献命令 (函数命令 → tex_command 金色) ------------------
引用表格第~\ref{tab:data}行，详见文献~\cite{key1,key2}。
\href{https://example.com}{链接文字}
% ↑ \ref, \cite, \href → tex_command 金色

% --- 12. 注释中的颜色「隔离」测试 ---------------------------------------
% 下面是关键测试：注释内的任何内容都应该是绿色，不能"泄露"其他颜色
% \begin{document}  ← 即使写了关键字命令，在注释中也应显示绿色
% $E=mc^2$          ← 数学模式在注释中应显示绿色
% \textbf{粗体}     ← 函数命令在注释中应显示绿色

\end{document}
% ← \end 是关键字命令 → tex_keyword 粉色
